\section[Small e: a bounded full-parent contact certificate for every r]{Small $e$: a bounded full-parent contact certificate for every $r$}

\begin{proposition}[Local part of the small-$e$ theorem]
Use the unchanged physical native C1 source, with $k,\zeta,e$ positive
infinitesimals and $\eta$ a positive bounded unit. Every native positive
$r=1-a\eta$, and every positive opening rate, is allowed.
At every independent test with $u,v,w$ infinitesimal, $Z=1$, and $w>0$,
the entire original parent is tensor SOS. Literal $\beta=0$ is included.
\end{proposition}
This proposition is the contact part only. The independent outer/state
join is supplied separately; no assertion here identifies a testing rate
with a source rate.

\paragraph{The source entry forced by the original sign.}
Keep $K,Q,E,\ell,\delta$ from (CA3), with $E=b^2$ and
$\ell=Qr=K-aE$. Here $E$ is a positive unit, $0<\ell<K$, and no lower
bound on $a$ or $\ell$ is assumed. Put
\[
 \varepsilon=1/\delta,\qquad j=\delta k,\qquad t=\varepsilon^2\tau.
\]
Both $\varepsilon,j$ are infinitesimal: $e=\varepsilon K/Q$ and
$\zeta=j^2QE/K$. The original condition $\beta\ge0$ gives
\[
 \frac{t}{e^2}\le\frac{2c_b}{\eta c_d},
\]
whose right side is bounded after cancelling the common native $r$
factor. Thus $\tau$ is positive bounded. The exact chart is
\[
 k=\varepsilon j,\quad K=1+\varepsilon^2j^2,\quad
 Q=K^2-\ell\varepsilon^2j^2,\quad a=(K-\ell)/E,
\]
\begin{align*}
 \eta&=KE/Q,\quad r=\ell/Q,\quad H=r+\eta,\quad
 e^2=\varepsilon^2K^2/Q^2,\\
 \zeta&=j^2QE/K,\quad R=K^2/Q,\quad T=K-k^2t\zeta. \tag{SE1}
\end{align*}
All denominators displayed here are positive units. The cancelled
coquotients $\widetilde c_i=c_i/\ell$ are the same short exact formulas
as in (CA3); the cancellation introduces no inverse of $\ell$.
Their residues are $2(E+\ell)$ for $i=2,b,3,5,\alpha$, and $4$ for
$i=d$; $c_4$ has residue $2$. These formulas hold uniformly as
$\ell$ tends to zero or to one.

\paragraph{One bounded centered remainder, including both normals.}
Set
\[
 A_* =\frac{ArH}{1+rH},\qquad
 A-A_* =\frac{A}{1+rH}>0.                              \tag{SE2}
\]
Retain the entire square $4(A-A_*)D_uD_v(Y-r uW)^2$.
It and its full first derivative vanish at exact native contact.
Divide the remaining parent by the positive pure coefficient $a$;
every factor $a$ cancels before coefficients are used. Write
\[
 \begin{array}{ll}
 b_2=8\widetilde c_2e^2/(R\widetilde c_3),&
 b_b=8\widetilde c_be^2/(R\widetilde c_3),\\
 b_3=8/T,&b_5=8\widetilde c_5/(T\widetilde c_3),\\
 L=Ar/a=2\eta^2c_4/(Q\widetilde c_3),&
 d_0=d/a=2\eta\widetilde c_d/\widetilde c_3,\\
 A_0=A_*/a=LH/(1+rH),&
 \alpha_0=\alpha/a=2\varepsilon^4j^6Q^2\eta^2
       \widetilde c_\alpha/(K^2\widetilde c_3).
 \end{array}                                                     \tag{SE3}
\]
Every quantity in (SE3) is bounded with unit denominators, and
$b_2,b_b$ contain $\varepsilon^2$.
Use the original normals $N,M_B$, put $p=t/R$, $h=a/T$, and write
\begin{align*}
 q_A&=vD_uw,\qquad q_B=v(D_u-uN-puD_uv w),\\
 W&=w+u(M_B+hq_B)/D_u,\qquad Y=uw-(M_B+hq_B)/D_u.
\end{align*}
The entire bounded remainder is
\begin{align*}
G_s={}&b_2\{N^2-2q_AN+(1-p)q_A^2\}
 +b_bD_v\{N^2+2q_AN+pq_A^2\}\\
 &+b_3D_v\{M_B^2+2q_BM_B+hq_B^2\}
 +b_5\{M_B^2-2q_BM_B+(1-h)q_B^2\}\\
 &+4A_0D_uD_v(Y-r uW)^2
 +4LD_uD_v(2uWY-r u^2W^2)\\
 &+4\alpha_0D_u(vY-1)^2
 +2d_0D_u\{(vY-1)^2-D_v\}
 +\nu D_vN^2,\qquad \nu=d_0(R-t)/R.                 \tag{SE4}
\end{align*}
This is (SD4), after the indicated exact centered subtraction and
normalization, except that a positive part of the original $\Pi_2$ is
kept inside the remainder. Indeed the original coefficient of
$D_vN^2$ in $\Pi_2/a$ is
\[
 \frac{4\beta+2d}{a}\frac{R-t}{R}\ \ge\ 2\nu.
\]
Thus subtracting only $\Pi_2/a-\nu D_vN^2$ remains legitimate, even
when $\beta=0$. The other three original packets are retained unchanged.
All full first normals and every hard quadratic term are present in
(SE4); there are no terms of higher normal degree.

\paragraph{A conductor that never divides by $\varepsilon$.}
Define the separated monomials and differences
\[
 A=\varepsilon wv,\quad B=v,\quad C=Euw,\quad
 D_0=\varepsilon^2j^3E,\quad k=\varepsilon j,
\]
\[
 x=B-C,\quad y=A-kC,\quad z=uB-kC,\quad h_0=kC-D_0.
\]
The two standard radical pairings are obtained by dividing the original
four squared circuit monomials by $\delta$: they are
$A,uB,kC,D_0$, and their product is $(k^2BC)^2$.
Their strict circuit is
$T_4=y^2+z^2+h_0^2-2yz+2q_1^2$.
The exact contact is $(u,w,v)=(k,j,\varepsilon j^2E)$.
Besides the three basic conductor identities,
\[
 C(u-k)=z-ux,\quad
 \varepsilon C(w-j)=y-\varepsilon wx,\quad
 C(v-\varepsilon j^2E)=E(wz+jy-kwx),                \tag{SE5}
\]
use the additional, division-free identity
\[
 C v(w-j)=Ew(w-j)z+Ejw y-Ekw^2x.                    \tag{SE6}
\]
For a monic contact division
$S=(u-k)T_u+(w-j)T_w+(v-\varepsilon j^2E)T_v$,
whenever $T_w=\varepsilon U+vV$, the bounded conductor is
\[
 \begin{aligned}
 CS={}&x[-uT_u-\varepsilon wU-Ekw^2V-EkwT_v]\\
 &+y[U+EjwV+EjT_v]
 +z[T_u+Ew(w-j)V+EwT_v].
 \end{aligned}                                      \tag{SE7}
\]
No division by $\varepsilon,j,k,r,a,t$, or a contact difference occurs.

Here is the full rational membership argument required for applying
(SE7), not only a finite truncation. At zero normals let
$V_0=1-puv w$. Then
\[
 q_A=D_u wv,\quad q_B=D_uvV_0,\quad
 uW=C/E+u^2hvV_0,\quad Y=C/E-hvV_0.
\]
The scalar of (SE4) therefore has an exact raw allocation
\[
 f_AA^2+f_{BB}B^2+f_{BC}BC+f_{CC}C^2+\rho D_0^2.     \tag{SE8}
\]
Expand only the nonconstant part of $(vY-1)^2$ into $B,C$; its constant
has $\alpha_0/D_0^2=2\widetilde c_\alpha/\widetilde c_3$.
The coefficients in (SE8) have unit denominators independent of $w$.
Every dependence on $w$ in them comes through $V_0$, and hence has a
factor $\varepsilon^2$; modulo $\varepsilon$ they are independent of
$w$ and even in $u,v$. Both complete first normals are in $(A,B,C)$
with bounded coefficients by the same exact formulas.

Clear the finitely many unit denominators, chosen even in $u,v$ and
independent of $w$, and use the unchanged scalar square-lift algorithm.
For its first scalar conductor, the only $w$-dependent coefficient
$T_w$ is divisible by $\varepsilon$ (in fact by $\varepsilon^2$).
At $\varepsilon=u=0$, its residual polynomial $S$ is even in $v$,
and contact makes its constant zero; thus $S\in(v^2)$ there.
More explicitly, modulo $\varepsilon$ the first divided differences
satisfy $T_u\in(u)$, $T_w=0$, and $T_v\in(v)$ after the preceding
monic substitutions. It follows
that after this first conductor, the sole possible $w$ dependence
modulo $\varepsilon,u$ in the four second-stage remainders is of the
form $w v$ times a polynomial in $v$. Hence their $T_w$ all belong to
$(\varepsilon,v)$, exactly the hypothesis of (SE7). The complete
normal raw coefficients are independent of $w$ modulo $\varepsilon$,
so their one conductor has $T_w\in(\varepsilon)$. The retained contact
squares and the original native balances supply the exact contact
value and first-jet identities throughout. This proves the full
scalar-square and complete first-normal lifts over the bounded ring.

\paragraph{Finite residues of that full lift.}
Give $(\varepsilon,j,u,w,v,N,M_B)$ weights $(1,1,2,1,3,5,3)$;
$E,\ell,\tau$ have weight zero. Every positive-weight variable is
infinitesimal, while all weight-zero coefficients are bounded.
The contact substitutions preserve the grading; $x$ has weight three
and the other circuit roots weight five. Both (SE5) and (SE6) preserve
weights, so after the scalar collection through weight ten and full
first normals through weight five, every lifted remainder coefficient
has positive weight. Restore the unit denominators afterward.

Set $H_0=E+\ell$ and $V_*=1+\ell H_0$. Direct exact multiplication of
(SE1)--(SE4), followed by (SE7), gives scalar residue
\[
       \frac{8(E+2)}{H_0V_*}x^2+8T_4,                         \tag{SE9}
\]
and complete first-normal residues
\[
       L_N=0,\qquad L_M=\frac{16E(H_0+1)}{H_0V_*}\,x.
\]
The hard block has residue
\[
 \operatorname{diag}\left(\frac{4E}{H_0},
           \frac{8(E^2+2E\ell+2\ell^2+2)}{V_*}\right).       \tag{SE10}
\]
These statements include coefficient corrections of positive weight;
they are not asserted as literal equalities of the full rational scalar.
The exact checker reconstructs (SE4), verifies the full finite lifts,
and compares their residues with (SE9)--(SE10).
The only hard-soft coupled determinant is
\[
 \det G_{M,x}
 =\frac{64(E^2+2E\ell+4E+4\ell)}{H_0^2V_*}>0.              \tag{SE11}
\]
Together with $4E/H_0>0$ and the strict two-pairing circuit this gives
the full bounded strict eight-root Gram. All original positive packets,
the centered square, and the positive pure-coefficient factor $a$ are
then restored. All inverses and radicals remain separated between the
two independent real closed fields.
